% ======================================================================
%  PART 12 -- Exam Strategy Checklist + Classic Pitfalls
% ======================================================================
\section{Exam Strategy Checklist}
\label{sec:strategy}

The exam (cf.\ the mock) is \textbf{closed-book, no calculator}, and made of a
few multi-part questions mixing \emph{stated definitions}, \emph{short proofs},
and \emph{calculations}. The recurring tasks are: write down a definition;
compute an ACVS/ACF; find an MA$(\infty)$ representation; check
causality/invertibility; derive a spectral density; forecast and give the
error; prove (joint) second-order stationarity; and read cross-spectra. This
checklist maps a \emph{question cue} to the \emph{tool} you deploy.

\begin{tcolorbox}[enhanced,breakable,colback=orange!6!white,colframe=orange!75!black,
  arc=1.5mm,boxrule=0.7pt,title=\textbf{Cue $\rightarrow$ Tool: the decision table},
  coltitle=white,colbacktitle=orange!75!black,fonttitle=\bfseries]
\renewcommand{\arraystretch}{1.35}
\begin{tabular}{@{}p{0.46\textwidth}@{\hspace{0.5em}}p{0.5\textwidth}@{}}
\toprule
\textbf{When the question says\dots} & \textbf{Deploy\dots}\\
\midrule
``State the definition of (weak/strict) stationarity'' &
The three moment conditions (mean const, finite const variance, covariance a
function of $\tau$ only) / equality of all finite-dim.\ distributions under
shift.\\
``Is $\{X_t\}$ (weakly) stationary?'' &
Check the 3 conditions in order; stop at first failure. Watch for undefined
moments and a residual $t$-dependence in the covariance.\\
``Compute the ACVS/ACF of an MA$(q)$'' &
$\acvs_\tau=\sigma^2\sum_j\theta_j\theta_{j+\abs\tau}$ for $\abs\tau\le q$, else
$0$. ACF \emph{cuts off} after lag $q$.\\
``Compute the ACVS of an AR/ARMA'' &
Multiply the defining equation by $X_{t-\tau}$, take expectations $\Rightarrow$
a \textbf{difference equation} $\acvs_\tau=\sum\phi_j\acvs_{\tau-j}$; solve with
initial conditions (handle small $\tau$ separately because of the MA/noise
terms).\\
``Find $\psi_j$ in $X_t=\sum\psi_j\eps_{t-j}$'' &
Write $X_t=\Theta(B)/\Phi(B)\,\eps_t$; expand by partial fractions /
$1/(1-gB)=\sum g^kB^k$, or match coefficients in $\Phi(B)\sum\psi_jB^j=\Theta(B)$.\\
``Is it causal/stationary? invertible?'' &
\textbf{Roots of $\Phi(z)$ outside the unit circle} $\Rightarrow$
causal/stationary; \textbf{roots of $\Theta(z)$ outside} $\Rightarrow$
invertible. (AR always invertible; MA always stationary.)\\
``Find the PACF'' &
Durbin--Levinson recursion on the ACF; for AR$(p)$, $\pacf_\tau=0$ for
$\tau>p$. PACF \emph{cuts off} after $p$.\\
``Identify a model from ACF/PACF plots'' &
ACF cuts at $q$ \& PACF tails $\Rightarrow$ MA$(q)$; PACF cuts at $p$ \& ACF
tails $\Rightarrow$ AR$(p)$; both tail $\Rightarrow$ ARMA; slow ACF decay
$\Rightarrow$ difference (non-stationary).\\
``Estimate AR parameters'' &
Yule--Walker $\hat\phi=\hat\Gamma^{-1}\hat\gamma$,
$\hat\sigma^2=\hat\gamma_0-\sum\hat\phi_j\hat\gamma_j$; or (forward/backward) LS.\\
``Estimate MA/ARMA by LS'' &
Reconstruct innovations recursively ($\eps_0=0$), minimise $\sum\hat\eps_t^2$.\\
``Remove a trend / make stationary'' &
Difference: $\diff^q$ kills a degree-$(q-1)$ polynomial trend; $\diff_{(s)}$
removes period-$s$ seasonality.\\
``Forecast $X_{n+h}$ (causal ARMA)'' &
$X^n_{n+h}=\sum_{j\ge h}\psi_j\eps_{n+h-j}$ (future $\eps\equiv0$); MSE
$=\sigma^2\sum_{j=0}^{h-1}\psi_j^2$. Or the recursive truncated recipe.\\
``Forecast an ARIMA'' &
Forecast the differenced ARMA $Z=\diff^dX$, then un-difference (cumulate).
$P^n_{n+h}\to\infty$.\\
``Give a prediction interval'' &
$X^n_{n+h}\pm z_{1-\alpha/2}\sqrt{\Var(e_n(h))}$ (Gaussian).\\
``Find the spectral density'' &
$\sdf(f)=\sum_\tau\acvs_\tau e^{-2\pi i f\tau}$; for ARMA use
$\sdf(f)=\sigma^2\abs{\Theta(e^{-2\pi i f})}^2/\abs{\Phi(e^{-2\pi i f})}^2$.\\
``Spectrum of a filtered series $Y=\filt X$'' &
$\sdf_Y(f)=\abs{H(f)}^2\sdf_X(f)$, $H(f)=\sum_m h_m e^{-2\pi i f m}$.\\
``Periodogram expectation/variance'' &
$\E[\hat\sdf^{(p)}]=\Fej_n*\sdf$ (biased, blurring); $\Var\to\sdf(f)^2$ (does
not vanish). Taper for bias, multitaper for variance.\\
``Is this $\acvs_\tau$ a valid ACVS?'' &
Compute its SDF; it is valid \emph{iff} $\sdf(f)\ge0$ for all $f$ (equivalently
$\acvs$ is positive semi-definite).\\
``Cross-spectrum / coherence of a delay'' &
$S_{1,2}(f)=c\,e^{-2\pi i f d}S_{2,2}(f)$: amplitude $\abs c\,S_{2,2}$, phase
$-2\pi f d$.\\
``Is the VAR stable?'' &
Roots of $\det(I-\Phi_1z-\cdots-\Phi_pz^p)$ outside unit circle (companion
eigenvalues inside).\\
``Test residual autocorrelation'' &
Ljung--Box $Q_m=n(n+2)\sum_{\tau=1}^m\hat r_\tau^2/(n-\tau)\sim\chi^2_{m-p-q}$.\\
``Compare models'' &
AIC/AICc/BIC ($k=p+q+1$); smaller is better; BIC penalises complexity more.\\
``Model changing variance / volatility'' &
ARCH/GARCH on the conditional variance; recall it is uncorrelated but
dependent; stationarity $\sum\alpha+\sum\beta<1$.\\
\bottomrule
\end{tabular}
\end{tcolorbox}

\begin{keytech}{The six moves you will almost certainly need}{p12-six}
\begin{enumerate}
  \item \textbf{Stationarity check} (3 conditions).
  \item \textbf{ACVS via the white-noise filter trick} (only equal noise
        indices survive) or via the \textbf{difference equation} (multiply by
        $X_{t-\tau}$, take expectations).
  \item \textbf{Root test} for causality/invertibility/stability.
  \item \textbf{$\Theta/\Phi$ expansion} for $\psi$-weights (partial fractions
        or coefficient matching).
  \item \textbf{Recursive forecasting} and its MSE.
  \item \textbf{Spectral density} from ACVS, from a filter, or from ARMA
        polynomials.
\end{enumerate}
\end{keytech}

\begin{keytech}{Exam hygiene}{p12-hygiene}
\begin{itemize}
  \item State which \emph{theorem/definition} you invoke (examiners reward
        naming the tool: ``by causality $\E[\eps_tX_{t-k}]=0$''; ``by Herglotz
        / the prediction equations / Isserlis'').
  \item No calculator: keep fractions exact; factor polynomials by hand
        ($z=\frac{-b\pm\sqrt{b^2-4ac}}{2a}$) and report $\abs{z}$ vs $1$.
  \item For ``$\abs z>1$?'' with complex roots, compare
        $\abs z^2=(\Re z)^2+(\Im z)^2$ to $1$ --- no need to take the square
        root.
  \item Always sanity-check: $\acvs_0\ge0$, $\abs{\acf_\tau}\le1$, MA(1) has
        $\abs{\acf_1}\le\tfrac12$, an SDF must be $\ge0$ and real, a spectral
        \emph{matrix} must be Hermitian.
\end{itemize}
\end{keytech}

% ======================================================================
\section{Classic Pitfalls to Avoid}
\label{sec:pitfalls}

\begin{pitfall}
\textbf{Roots inside vs outside the unit circle.} Stationarity/causality
requires the roots of $\Phi(z)$ to lie \emph{outside} $\abs z=1$ (equivalently
$\Phi(z)\ne0$ for $\abs z\le1$). Students routinely flip this. The
\emph{characteristic} root condition (``$\abs z>1$'') is the opposite of the
companion-eigenvalue condition (``$\abs\lambda<1$''): both say the same thing,
so do not mix the two phrasings in one argument.
\end{pitfall}

\begin{pitfall}
\textbf{Strict vs weak stationarity.} They do not imply each other. Strict but
not weak: i.i.d.\ Cauchy (no finite variance). Weak but not strict: a process
with matched first two moments but non-Gaussian higher structure. They coincide
\emph{only} for Gaussian processes.
\end{pitfall}

\begin{pitfall}
\textbf{Lag goes in the FIRST argument (multivariate).}
$\Gamma_\tau=\Cov(X_{t+\tau},X_t)$, so $\Gamma_{-\tau}=\Gamma_\tau\Tr$ (not
$\Gamma_\tau$). Cross-covariance sequences are \emph{not} symmetric in $\tau$:
$\gamma^{(j,k)}_\tau=\gamma^{(k,j)}_{-\tau}$. In the univariate real case this
distinction is invisible, which is exactly why it trips people in the
multivariate case.
\end{pitfall}

\begin{pitfall}
\textbf{ACF cut-off, PACF cut-off --- which is which.} MA$(q)$: \emph{ACF} cuts
off after $q$ (PACF tails). AR$(p)$: \emph{PACF} cuts off after $p$ (ACF tails).
ARMA: both tail off. Swapping these is the single most common identification
error.
\end{pitfall}

\begin{pitfall}
\textbf{Identifiability of MA / sign of $\theta_0$.} Fixing $\theta_0=-1$ is not
cosmetic: without it, scaling
$(\theta_j,\sigma^2)\mapsto(c\theta_j,c^{-2}\sigma^2)$ gives the \emph{same}
ACVS, so the model is unidentifiable. Likewise a non-invertible MA and its
invertible ``twin'' (root $g$ vs $1/g$) share the same ACF --- pick the
invertible one.
\end{pitfall}

\begin{pitfall}
\textbf{Biased is not bad.} The divide-by-$n$ ACVS estimator $\hat\acvs_\tau$ is
biased yet positive semi-definite and has \emph{smaller} MSE than the unbiased
$\tilde\acvs_\tau$. Unbiasedness is not the objective; positive
semi-definiteness and low MSE are. Do not ``correct'' the estimator on autopilot.
\end{pitfall}

\begin{pitfall}
\textbf{The periodogram is not consistent.} Its variance tends to $\sdf(f)^2$
and does \emph{not} shrink with $n$ --- more data makes it wigglier, not
smoother. Reducing bias needs \emph{tapering}; reducing variance needs
\emph{smoothing/multitapering} ($\Var\approx\sdf(f)^2/K$).
\end{pitfall}

\begin{pitfall}
\textbf{Forecasts of stationary vs integrated processes.} A stationary ARMA
forecast reverts to the mean and its error variance saturates at $\gamma_0$. An
\emph{integrated} (ARIMA, $d\ge1$) forecast does \emph{not} revert and its error
variance $\to\infty$. Quoting a finite long-horizon variance for a random walk
is wrong.
\end{pitfall}

\begin{pitfall}
\textbf{Setting future innovations to zero (only).} In the $\psi$-weight
forecast, set $\eps_{n+1},\dots,\eps_{n+h}=0$ but keep the \emph{observed past}
innovations as residuals; in the recursive recipe, replace future
\emph{observations} by their forecasts, future errors by $0$, and past errors by
residuals --- do not zero the past errors.
\end{pitfall}

\begin{pitfall}
\textbf{ARCH/GARCH are uncorrelated but dependent.} $\Corr(r_{t+\tau},r_t)=0$,
yet $\Corr(r^2_{t+\tau},r^2_t)\ne0$ (geometric for ARCH(1):
$\alpha_1^{\abs\tau}$). ``Uncorrelated'' $\ne$ ``independent'' here. Also check
\emph{moment} conditions ($\alpha_1<1$ for variance, $\alpha_1^2<1/3$ for finite
kurtosis under Gaussian noise) separately from stationarity.
\end{pitfall}

\begin{pitfall}
\textbf{Ljung--Box degrees of freedom.} Use $\chi^2_{m-p-q}$, subtracting the
number of fitted ARMA parameters --- not $\chi^2_m$. Forgetting the $-(p+q)$
shifts the critical value and can flip the conclusion.
\end{pitfall}

\begin{pitfall}
\textbf{Angular vs ordinary frequency.} This course uses ordinary frequency $f$
(period $1/f$), so the SDF integrates over $[-\tfrac12,\tfrac12]$ and
$\sdf(f)=\sum\acvs_\tau e^{-2\pi i f\tau}$. Switching to $\omega=2\pi f$ scatters
stray $2\pi$ factors --- keep one convention.
\end{pitfall}

\begin{pitfall}
\textbf{A valid ACVS must be positive semi-definite.} A plausible-looking
sequence (e.g.\ $\acvs_\tau=1$ for $\abs\tau\le K$, else $0$) need \emph{not} be
a real ACVS: its SDF can go negative. Always test via $\sdf(f)\ge0$ before
claiming a sequence is an autocovariance.
\end{pitfall}

\begin{center}
\begin{tcolorbox}[width=0.9\textwidth,colback=green!5!white,colframe=green!48!black,
  arc=2mm,boxrule=0.8pt,halign=center,fontupper=\large\bfseries]
You are ready when you can do every \textcolor{cyan!55!black}{Exercise} in this
document unaided, name the tool behind every \textcolor{red!62!black}{Theorem},
and avoid every \textcolor{red!70!black}{Pitfall} above. Good luck!
\end{tcolorbox}
\end{center}
